\documentclass[10pt,a4paper]{amsart}
\usepackage[margin=19mm]{geometry}
\usepackage{amsmath,amssymb,mathtools}
\usepackage[scr=rsfs,cal=euler]{mathalfa}
\usepackage{xfrac}
\usepackage[unicode,hidelinks,bookmarks=false]{hyperref}
\newcommand{\ie}{i.e.\ }
\newcommand{\cf}{cf.\ }
\newcommand{\resp}{respectively\ }
\newcommand{\Subsection}{\subsection}
\providecommand{\nobreakdash}{\nobreak\hbox{-}}
\setlength{\emergencystretch}{2em}
\multlinegap0pt
\usepackage{amssymb}
\usepackage{mathtools}
\usepackage{bm}
\usepackage[shortlabels]{enumitem}
\usepackage{xcolor}
\newtheorem{assumption}{Assumption}
\newtheorem{theorem}{Theorem}
\newtheorem{proposition}{Proposition}
\newtheorem{lemma}{Lemma}
\newcommand{\Coll}{\mathscr C}
\def\E{\mathbb{E}}
\newcommand{\Ent}{\operatorname{Ent}}
\def\Fc{{\mathcal F}}
\def\P{\mathbb{P}}
\def\R{\mathbb{R}}
\def\Xc{{\mathcal X}}
\newcommand{\dd}{\mathrm{d}}
\def\hX{\widehat{\Xc}^{\,N}}
\def\id{\mathbf{1}}
\newcommand{\pmax}{p_{\max}}
\newcommand{\pmin}{p_{\min}}
\newcommand{\tr}{\operatorname{tr}}
\title{Semi-discrete quadratic Wasserstein energy and state-dependent Langevin exploration}
\author{Ran Gu, Gaoyue Guo, Kelvin Shuangjian Zhang}
\date{Source: arXiv:2609.03405v1 (CC BY 4.0)}
\begin{document}
\maketitle

\noindent\emph{Source and licence.} Semi-discrete quadratic Wasserstein energy and state-dependent Langevin exploration. Version arXiv:2609.03405v1 (CC BY 4.0), distributed by arXiv under the Creative Commons Attribution 4.0 International licence, http://creativecommons.org/licenses/by/4.0/, as recorded in the arXiv licence metadata for this exact version and shown on the abstract page https://arxiv.org/abs/2609.03405v1. Full licence notice: \texttt{licenses/arxiv-2609.03405-CC-BY-4.0.txt}.\par\smallskip

\noindent\textit{Editorial scope.} Counted unit: the article's theorem thm:HJB (source0.tex lines 542--605). The article's complete original proof of it is delivered with it (source0.tex lines 2837--3049, one \texttt{proof} environment of 213 lines, which the article places away from the statement). No mathematics is written, repaired or re-derived: the statement and the proof are the article's own lines, in the article's own notation, and no retained line is substituted or re-flowed.  Retained uncounted as the source-local context the proof needs: lines 219--245; lines 533--540; lines 1851--1853; lines 1865--1868; lines 1901--1908; lines 1932--1936; lines 1963--1986; lines 1979--1984; lines 2212--2228; lines 2378--2391; lines 2380--2384; lines 2636--2660; lines 2651--2657; lines 2738--2744; lines 2746--2763; lines 2812--2821.  Nothing was omitted from any retained span, so the package carries no omission record.  No retained line contains a citation, so the wrapper carries no bibliography and no bibliography entry.  No retained line contains a figure environment, an image-inclusion command or drawing code, so no image is delivered and the package declares no asset dependency.  Editorial formatting only, in addition: an AMS \texttt{amsart} preamble that restates the article's own declarations and macros used by the retained lines, namely the environments \texttt{assumption}, \texttt{theorem}, \texttt{proposition}, \texttt{lemma} on separate counters, together with the standard packages those lines need.  The retained environments print in this document as Assumption 1, Theorem 1, Proposition 1, Lemma 1 in order of appearance, whereas the article prints the same items with the numbering of its own full document.  The complete native source of the article as distributed in the arXiv e-print for this exact version is bundled unchanged as \texttt{sources/209295/source0.tex}.
\begin{assumption}[Target measure and prescribed masses]\label{ass:geometry}
{\rm (i)} The set $\Omega\subset\R^d$ is compact and convex with nonempty interior $\Omega^\circ$.

\medskip

\noindent {\rm (ii)} The probability measure $\nu$ has the form
\begin{align}
 \nu(\dd y)&=\rho(y)\id_{\Omega}(y)\dd y,
 \label{eq:density-assumption-main}\\
 \rho&\in C(\Omega)\cap W^{1,1}(\Omega^\circ),
 \qquad
 0<\underline\rho:=\min_{y\in\Omega}\rho(y)
 \le \max_{y\in\Omega}\rho(y)=:\overline\rho<\infty,
 \notag\\
 \int_\Omega\rho(y)\dd y&=1.
 \notag
\end{align}

\noindent {\rm (iii)} The prescribed masses $p_1,\ldots,p_N$ satisfy
\begin{equation}\label{eq:weights-main}
 \sum_{i=1}^Np_i=1,
 \qquad
 0<\pmin:=\min_{1\le i\le N}p_i,
 \qquad
 \pmax:=\max_{1\le i\le N}p_i.
\end{equation}
\end{assumption}

\begin{equation}\label{eq:exploratory-HJB-main-results}
 -\vartheta v_\lambda
 -\nabla E\cdot\nabla v_\lambda
 +E
 +\mathscr H_\lambda(\Delta v_\lambda)
 =0
 \qquad\text{on }\hX.
\end{equation}

\begin{equation}\label{eq:admissible-control-class}
 \pi:[0,\infty)\times\Omega_W\longrightarrow\mathfrak U
\end{equation}

\begin{equation}\label{eq:controlled-SDE}
 \dd X_t^\pi=-\nabla E(X_t^\pi)\dd t+\sigma_t^\pi\dd W_t,
 \qquad X_0^\pi=X.
\end{equation}

\begin{align}
 \mathscr H_\lambda'(q)
 &=\int_a^1u\,\varpi_q^\lambda(\dd u)\in(a,1),
 \label{eq:H-prime}\\
 \mathscr H_\lambda''(q)
 &=-\frac1\lambda\operatorname{Var}_{\varpi_q^\lambda}(u)<0.
 \label{eq:H-second}
\end{align}

\begin{equation}\label{eq:gibbs-variational}
 \inf_{\varpi\in\mathfrak U}
 \left\{q\bar u(\varpi)+\lambda\Ent_{[a,1]}(\varpi)\right\}
 =\mathscr H_\lambda(q),
\end{equation}

\begin{proposition}\label{prop:controlled-dynamics}
For every relaxed control $\pi$ and every initial condition $X\in\hX$, there exists a unique
continuous process $X^\pi=(X_t^\pi)_{t\ge0}$, adapted to the raw canonical filtration
$(\Fc_t^0)_{t\ge0}$, such that, almost surely,
\begin{equation}\label{eq:controlled-integral-form}
 X_t^\pi
 =X-\int_0^t\nabla E(X_s^\pi)\dd s
 +\int_0^t\sigma_s^\pi\dd W_s,
 \qquad t\ge0.
\end{equation}
In other words, \eqref{eq:controlled-SDE} has a pathwise unique global strong solution.  This
solution never reaches the collision set:
\begin{equation}\label{eq:no-collision-control}
 \P_X\left[X_t^\pi\in\hX\text{ for every }t\ge0\right]=1.
\end{equation}
It also satisfies
\begin{equation}\label{eq:moment-controlled}
 \E_X|X_t^\pi|^2
 \le e^{-\pmin t}|X|^2+
 \frac{m_2(\nu)+2m}{\pmin},
 \qquad t\ge0,
\end{equation}
uniformly over all relaxed controls.
\end{proposition}

\begin{equation}
 \E_X|X_t^\pi|^2
 \le e^{-\pmin t}|X|^2+
 \frac{m_2(\nu)+2m}{\pmin},
 \qquad t\ge0,
\end{equation}

\begin{lemma}\label{lem:uniform-localization-control}
Let $K\Subset\hX$ and $T>0$.  Uniformly over all relaxed controls and all initial states $X\in K$,
\begin{equation}\label{eq:uniform-fourth-moment-control}
 \sup_{\pi}\sup_{X\in K}
 \E_X\left[\sup_{0\le t\le T}|X_t^\pi|^4\right]<\infty.
\end{equation}
Moreover,
\begin{align}
 \lim_{R\to\infty}\sup_{\pi}\sup_{X\in K}
 \P_X\left(\sup_{0\le t\le T}|X_t^\pi|\ge R\right)&=0,
 \label{eq:uniform-no-escape-control}\\
 \lim_{\delta\downarrow0}\sup_{\pi}\sup_{X\in K}
 \P_X\left(\min_{i<j}\inf_{0\le t\le T}
 |X_t^{\pi,i}-X_t^{\pi,j}|\le\delta\right)&=0.
 \label{eq:uniform-no-near-collision-control}
\end{align}
\end{lemma}

\begin{lemma}\label{lem:value-continuity}
The value function is finite and satisfies
\begin{equation}\label{eq:value-growth}
 -\frac{\lambda\log(1-a)}{\vartheta}
 \le v_\lambda(X)\le C_\lambda(1+|X|^2),
 \qquad X\in\hX.
\end{equation}
Moreover, for every compact set $K\Subset\hX$, the family of maps
\[
 X\longmapsto J_\lambda(X;\pi),
\]
indexed by controls having finite cost at one point of $K$, is equicontinuous on $K$.  Every such
control has finite cost at all points of $K$.  Consequently, $v_\lambda$ is continuous on $\hX$.
\end{lemma}

\begin{equation}
 -\frac{\lambda\log(1-a)}{\vartheta}
 \le v_\lambda(X)\le C_\lambda(1+|X|^2),
 \qquad X\in\hX.
\end{equation}

\begin{proposition}\label{prop:DPP}
Let $\tau$ be any bounded stopping time for the raw canonical filtration $\mathbb F^0$.  Then, for
every $X\in\hX$,
\begin{equation}\label{eq:DPP-HJB-detailed}
 v_\lambda(X)=\inf_\pi\E_X\left[
 \int_0^\tau e^{-\vartheta s}\ell_\lambda(X_s^\pi,\pi_s)\dd s
 +e^{-\vartheta\tau}v_\lambda(X_\tau^\pi)\right].
\end{equation}
The same principle holds for control-dependent localized exit times.  Namely, if
$D\Subset\hX$ is open, $T>0$, and
\begin{equation}\label{eq:localized-stopping-time-DPP}
 \tau_{D,T}^\pi
 :=T\wedge\inf\{s\ge0:X_s^\pi\notin D\},
\end{equation}
then, for every $X\in D$,
\begin{equation}\label{eq:DPP-stopping-control}
 v_\lambda(X)=\inf_\pi\E_X\left[
 \int_0^{\tau_{D,T}^\pi}e^{-\vartheta s}
 \ell_\lambda(X_s^\pi,\pi_s)\dd s
 +e^{-\vartheta\tau_{D,T}^\pi}
 v_\lambda(X_{\tau_{D,T}^\pi}^\pi)\right].
\end{equation}
In both formulas, controls for which the displayed expectation is infinite may be omitted from the
infimum.
\end{proposition}

\begin{equation}
 v_\lambda(X)=\inf_\pi\E_X\left[
 \int_0^{\tau_{D,T}^\pi}e^{-\vartheta s}
 \ell_\lambda(X_s^\pi,\pi_s)\dd s
 +e^{-\vartheta\tau_{D,T}^\pi}
 v_\lambda(X_{\tau_{D,T}^\pi}^\pi)\right].
\end{equation}

\begin{equation}\label{eq:exploratory-HJB}
 -\vartheta v_\lambda(X)
 -\nabla E(X)\cdot\nabla v_\lambda(X)
 +E(X)
 +\mathscr H_\lambda\bigl(\Delta v_\lambda(X)\bigr)=0,
 \qquad X\in\hX.
\end{equation}

\begin{lemma}\label{lem:HJB-interior-regularity}
Let $D'\Subset D\Subset\hX$, and let $u$ be a bounded viscosity solution on $D$ of
\[
 -\vartheta u-\nabla E\cdot\nabla u+E+\mathscr H_\lambda(\Delta u)=0.
\]
There exists
$\bar\beta=\bar\beta(m,a)\in(0,1)$ such that, for every
$0<\beta<\bar\beta$,
\[
 u\in C^{2,\beta}(D').
\]
The corresponding interior norm is bounded in terms of
\[
 \operatorname{dist}(D',\partial D),\quad m,\quad a,\quad
 \vartheta,\quad\lambda,\quad
 \|E\|_{C^2(D)},\quad\sup_D|u|.
\]
\end{lemma}

\begin{lemma}\label{lem:laplacian-C1}
Let $u\in C_{\mathrm{loc}}^{2,\beta}(\hX)$ be a viscosity solution of
\[
 -\vartheta u-\nabla E\cdot\nabla u+E+\mathscr H_\lambda(\Delta u)=0.
\]
Then the equation holds pointwise and
\begin{equation}\label{eq:laplacian-C1-lemma}
 \Delta u\in C_{\mathrm{loc}}^1(\hX).
\end{equation}
\end{lemma}

\begin{theorem}\label{thm:HJB}
Suppose that Assumption~\ref{ass:geometry} holds and that $d\ge2$.  Then $v_\lambda$ is finite and
continuous on $\hX$, has at most quadratic growth, and is a viscosity solution of
\eqref{eq:exploratory-HJB-main-results}.  There exists $\beta\in(0,1)$ such that
\begin{equation}\label{eq:v-C2beta}
 v_\lambda\in C^{2,\beta}_{\mathrm{loc}}(\hX).
\end{equation}
Moreover,
\begin{equation}\label{eq:q-C1}
 q_\lambda:=\Delta v_\lambda\in C^1_{\mathrm{loc}}(\hX),
\end{equation}
so the HJB equation holds classically on $\hX$.

For every $X\in\hX$, the pointwise Hamiltonian problem
\begin{equation}\label{eq:pointwise-Hamiltonian-minimization}
 \inf_{\varpi\in\mathcal P([a,1])}
 \left\{
 \Delta v_\lambda(X)\,\bar u(\varpi)
 +\lambda\Ent_{[a,1]}(\varpi)
 \right\}
\end{equation}
has a unique minimizer $\varpi_\lambda^*(X)$.  It has density
\begin{equation}\label{eq:optimal-density}
 \pi_\lambda^*(u;X)
 :=
 \frac{\exp[-u\Delta v_\lambda(X)/\lambda]}
 {\displaystyle\int_a^1\exp[-s\Delta v_\lambda(X)/\lambda]\dd s},
 \qquad u\in[a,1],
\end{equation}
that is,
$\varpi_\lambda^*(X)(\dd u)=\pi_\lambda^*(u;X)\dd u$.  Its mean temperature is
\begin{equation}\label{eq:optimal-mean-temperature}
 \tau_\lambda(X)
 :=\bar u\bigl(\varpi_\lambda^*(X)\bigr)
 =\int_a^1u\pi_\lambda^*(u;X)\dd u
 =\mathscr H_\lambda'\bigl(\Delta v_\lambda(X)\bigr)
 \in(a,1).
\end{equation}
The diffusion coefficient
\begin{equation}\label{eq:h-feedback}
 h_\lambda(X):=\sqrt{2\tau_\lambda(X)}
\end{equation}
belongs to $C^1_{\mathrm{loc}}(\hX)$ and satisfies
\[
 \sqrt{2a}\le h_\lambda(X)\le\sqrt2,
 \qquad X\in\hX.
\]
For every $X\in\hX$, the feedback equation
\begin{equation}\label{eq:optimal-feedback-SDE}
 \dd X_t^*
 =-\nabla E(X_t^*)\dd t+h_\lambda(X_t^*)\dd W_t,
 \qquad X_0^*=X,
\end{equation}
has a unique global strong solution which almost surely never reaches $\Coll$.  Finally,
\begin{equation}\label{eq:optimal-relaxed-control-process}
 \pi_t^*(\dd u)
 :=\varpi_\lambda^*(X_t^*)(\dd u)
 =\pi_\lambda^*(u;X_t^*)\dd u
\end{equation}
has finite cost and is optimal:
\begin{equation}\label{eq:verification-value}
 v_\lambda(X)=J_\lambda(X;\pi^*)=\inf_\pi J_\lambda(X;\pi).
\end{equation}
\end{theorem}

\begin{proof}[Proof of Theorem~\ref{thm:HJB}]
We indicate explicitly how the preceding auxiliary results enter the proof.

\medskip
\noindent\textit{Step 1: finiteness, growth, and continuity.}
These conclusions are exactly Lemma~\ref{lem:value-continuity}.  The proof of that lemma uses the
uniform localization result of Lemma~\ref{lem:uniform-localization-control}, but not the dynamic
programming principle.

\medskip
\noindent\textit{Step 2: viscosity solution property.}
Define the proper operator
\begin{equation}\label{eq:proper-HJB-operator}
 \Fc_\lambda(X,r,p,A)
 :=\vartheta r+\nabla E(X)\cdot p-E(X)
 -\mathscr H_\lambda(\tr A).
\end{equation}
Equation \eqref{eq:exploratory-HJB} is equivalent to
$\Fc_\lambda(X,v,\nabla v,D^2v)=0$.
Fix $X_0\in\hX$ and choose $r>0$ such that
$\overline B_{2r}(X_0)\Subset\hX$.  For a control $\pi$ and $h>0$, set
\begin{equation}\label{eq:viscosity-exit-time-control}
 \tau_h^\pi
 :=h\wedge\inf\{t\ge0:X_t^\pi\notin B_r(X_0)\}.
\end{equation}
On $B_r(X_0)$ the drift is bounded, uniformly over controls, and the diffusion coefficient is
bounded by $\sqrt2$.  The stopped integral equation and the Burkholder--Davis--Gundy inequality
therefore give
\begin{equation}\label{eq:short-time-state-estimate}
 \sup_\pi\E_{X_0}\left[
 \sup_{0\le s\le\tau_h^\pi}|X_s^\pi-X_0|^2\right]
 \le C_r(h+h^2).
\end{equation}
Indeed, the drift contribution is bounded by $C_rh^2$, while the quadratic expectation of the
stopped stochastic integral is bounded by $2mh$.  In particular,
\begin{equation}\label{eq:short-time-exit-probability-control}
 \sup_\pi\P_{X_0}(\tau_h^\pi<h)
 \le r^{-2}C_r(h+h^2)\longrightarrow0.
\end{equation}
Consequently, $\tau_h^\pi/h\to1$ in probability uniformly over controls.  If $G$ is continuous on
$\overline B_r(X_0)$, uniform continuity of $G$, \eqref{eq:short-time-state-estimate}, and
\eqref{eq:short-time-exit-probability-control} yield
\begin{equation}\label{eq:short-time-average-control}
 \frac1h\E_{X_0}\int_0^{\tau_h^\pi}e^{-\vartheta s}G(X_s^\pi)\dd s
 \longrightarrow G(X_0),
\end{equation}
uniformly over controls whenever $G$ itself does not depend on the chosen control.
Suppose first that $v_\lambda-\varphi$ has a local maximum zero at $X_0$, with
$\varphi\in C^2(B_{2r}(X_0))$.  After decreasing $r$ if necessary,
$v_\lambda\le\varphi$ on $\overline B_r(X_0)$.  Fix a constant action
$\varpi\in\mathfrak U$ with finite entropy and apply the stopping-time DPP
\eqref{eq:DPP-stopping-control} with $D=B_r(X_0)$ and horizon $h$.  Since the value is the infimum,
using this particular constant control gives
\[
 0\le\E_{X_0}\left[
 \int_0^{\tau_h^\varpi}e^{-\vartheta s}
 \ell_\lambda(X_s^\varpi,\varpi)\dd s
 +e^{-\vartheta\tau_h^\varpi}\varphi(X_{\tau_h^\varpi}^\varpi)
 -\varphi(X_0)\right].
\]
It\^o's formula for $e^{-\vartheta t}\varphi(X_t^\varpi)$ up to
$\tau_h^\varpi$ turns this into
\begin{align*}
 0\le\E_{X_0}\int_0^{\tau_h^\varpi}e^{-\vartheta s}
 \Bigl[&E(X_s^\varpi)+\lambda\Ent_{[a,1]}(\varpi)
 -\vartheta\varphi(X_s^\varpi)\\
 &-\nabla E(X_s^\varpi)\cdot\nabla\varphi(X_s^\varpi)
 +\bar u(\varpi)\Delta\varphi(X_s^\varpi)\Bigr]\dd s.
\end{align*}
Divide by $h$ and let $h\downarrow0$ using
\eqref{eq:short-time-average-control}.  Since $\varpi$ is arbitrary,
\eqref{eq:gibbs-variational} gives
\[
 E(X_0)-\vartheta\varphi(X_0)
 -\nabla E(X_0)\cdot\nabla\varphi(X_0)
 +\mathscr H_\lambda(\Delta\varphi(X_0))\ge0,
\]
which is equivalent to
\[
 \Fc_\lambda(X_0,\varphi(X_0),\nabla\varphi(X_0),D^2\varphi(X_0))\le0.
\]
Thus $v_\lambda$ is a viscosity subsolution.
Suppose now that $v_\lambda-\varphi$ has a local minimum zero at $X_0$.  Then
$v_\lambda\ge\varphi$ on $\overline B_r(X_0)$.  In the stopping-time DPP choose a control
$\pi^h$ whose displayed value is within $\varepsilon_h:=h^2$ of the infimum.  The terminal
comparison with $\varphi$ and It\^o's formula give
\begin{align*}
 h^2\ge\E_{X_0}\int_0^{\tau_h^{\pi^h}}e^{-\vartheta s}
 \Bigl[&E(X_s^{\pi^h})+\lambda\Ent_{[a,1]}(\pi_s^h)
 -\vartheta\varphi(X_s^{\pi^h})\\
 &-\nabla E(X_s^{\pi^h})\cdot\nabla\varphi(X_s^{\pi^h})
 +\bar u(\pi_s^h)\Delta\varphi(X_s^{\pi^h})\Bigr]\dd s.
\end{align*}
Pointwise in $(s,\omega)$, the Gibbs inequality yields
\[
 \bar u(\pi_s^h)\Delta\varphi(X_s^{\pi^h})
 +\lambda\Ent_{[a,1]}(\pi_s^h)
 \ge\mathscr H_\lambda(\Delta\varphi(X_s^{\pi^h})).
\]
The remaining integrand is now a continuous function of the state only.  Divide by $h$ and use
\eqref{eq:short-time-average-control}; since $h^2/h\to0$, we obtain
\[
 \Fc_\lambda(X_0,\varphi(X_0),\nabla\varphi(X_0),D^2\varphi(X_0))\ge0.
\]
This proves that $v_\lambda$ is a viscosity solution.  The role of
Proposition~\ref{prop:DPP} in the main theorem is precisely this step.

\medskip
\noindent\textit{Step 3: interior regularity and differentiability of the Laplacian.}
Lemma~\ref{lem:HJB-interior-regularity}, applied on arbitrary domains
$D'\Subset D\Subset\hX$, yields a number $\beta\in(0,1)$ such that
\[
 v_\lambda\in C_{\mathrm{loc}}^{2,\beta}(\hX).
\]
This proves \eqref{eq:v-C2beta}.  Lemma~\ref{lem:laplacian-C1} then gives
\[
 q_\lambda:=\Delta v_\lambda\in C_{\mathrm{loc}}^1(\hX),
\]
which is \eqref{eq:q-C1}.  In particular, the HJB equation holds classically on $\hX$.

\medskip
\noindent\textit{Step 4: construction and admissibility of the feedback.}
For $X\in\hX$, the unique minimizer of the pointwise Hamiltonian is
\[
 \varpi_\lambda^*(X)(\dd u)
 :=\varpi_{q_\lambda(X)}^\lambda(\dd u)
 =\pi_\lambda^*(u;X)\dd u,
\]
where $\pi_\lambda^*$ is given by \eqref{eq:optimal-density}.  Equations
\eqref{eq:H-prime} and \eqref{eq:q-C1} imply
\[
 \tau_\lambda,h_\lambda\in C^1_{\mathrm{loc}}(\hX),
 \qquad
 \sqrt{2a}\le h_\lambda\le\sqrt2.
\]
Thus both coefficients of the feedback equation \eqref{eq:optimal-feedback-SDE} are locally
Lipschitz on $\hX$.  The Picard and exhaustion construction in Step~1 of
Proposition~\ref{prop:controlled-dynamics}, now applied to the two state-dependent coefficients,
produces a pathwise unique maximal strong solution which can be chosen continuous and adapted to
the raw canonical filtration.  The Lyapunov calculation of Step~2 there is unchanged because
$h_\lambda^2=2\tau_\lambda\le2$.  In the localized collision argument, Girsanov removes the drift
and each pairwise difference has quadratic clock
\[
 4\int_0^t\tau_\lambda(X_s^*)\dd s,
\]
which lies between $4at$ and $4t$.  The same polarity argument therefore excludes collisions, and
the feedback solution is global.  The same Lyapunov calculation also yields the estimate
\eqref{eq:moment-controlled} with $X^\pi$ replaced by $X^*$.
The map $X\mapsto\varpi_\lambda^*(X)$ is continuous from $\hX$ to $\mathfrak U$.  Since $X^*$ is
continuous and raw adapted, it is raw progressive; hence
\begin{equation}\label{eq:feedback-control-raw-progressive}
 \pi_t^*(\dd u)
 :=\varpi_\lambda^*(X_t^*)(\dd u)
 =\pi_\lambda^*(u;X_t^*)\dd u
\end{equation}
is raw progressively measurable and is admissible in the sense of
\eqref{eq:admissible-control-class}.

\medskip
\noindent\textit{Step 5: verification and finite feedback cost.}
Let $\pi$ be any finite-cost control and set
$q_t:=\Delta v_\lambda(X_t^\pi)$.  The Gibbs inequality gives
\begin{equation}\label{eq:gibbs-verification-detailed}
 \bar u(\pi_t)q_t+\lambda\Ent_{[a,1]}(\pi_t)
 \ge\mathscr H_\lambda(q_t),
\end{equation}
with equality if and only if
$\pi_t(\dd u)=\pi_\lambda^*(u;X_t^\pi)\dd u$.
Choose the exhaustion
$D_k=\{Y\in\hX:|Y|<k,\ \mathfrak d(Y)>k^{-1}\}$ and let
$\sigma_k$ be the first exit time of $X^\pi$ from $D_k$.  Since
$v_\lambda$, $\nabla v_\lambda$, and $D^2v_\lambda$ are bounded on
$\overline D_k$, It\^o's formula for
$e^{-\vartheta t}v_\lambda(X_t^\pi)$ up to $T\wedge\sigma_k$ has a martingale with zero expectation.
Using the classical HJB equation and \eqref{eq:gibbs-verification-detailed} yields
\begin{align}
 &\E_X\left[e^{-\vartheta(T\wedge\sigma_k)}
 v_\lambda(X_{T\wedge\sigma_k}^\pi)\right]-v_\lambda(X)\notag\\
 &\qquad\ge-\E_X\int_0^{T\wedge\sigma_k}e^{-\vartheta t}
 \ell_\lambda(X_t^\pi,\pi_t)\dd t.
 \label{eq:Ito-verification-detailed}
\end{align}
By Lemma~\ref{lem:uniform-localization-control},
$\sup_{s\le T}|X_s^\pi|^4$ is integrable.  Together with the quadratic growth
\eqref{eq:value-growth}, this makes
$v_\lambda(X_{T\wedge\sigma_k}^\pi)$ uniformly integrable in $k$.  Since
$\sigma_k\to\infty$ almost surely and the running cost is nonnegative, we may let $k\to\infty$ in
\eqref{eq:Ito-verification-detailed}.  Furthermore, \eqref{eq:value-growth} and
\eqref{eq:moment-controlled} give
\begin{equation}\label{eq:terminal-vanishes-verification}
 \E_X\left[e^{-\vartheta T}|v_\lambda(X_T^\pi)|\right]
 \le C e^{-\vartheta T}
 \left(1+e^{-\pmin T}|X|^2+\frac{m_2(\nu)+2m}{\pmin}\right)
 \longrightarrow0.
\end{equation}
Letting $T\to\infty$ proves
$v_\lambda(X)\le J_\lambda(X;\pi)$.  For an infinite-cost control this inequality is automatic.
For the feedback process $X^*$, equality holds in
\eqref{eq:gibbs-verification-detailed}.  The stopped identity becomes
\[
 \E_X\int_0^{T\wedge\sigma_k}e^{-\vartheta t}
 \ell_\lambda(X_t^*,\pi_t^*)\dd t
 =v_\lambda(X)-
 \E_X\left[e^{-\vartheta(T\wedge\sigma_k)}
 v_\lambda(X_{T\wedge\sigma_k}^*)\right]
 \le v_\lambda(X),
\]
because \eqref{eq:value-growth} implies
$v_\lambda\ge-\lambda\log(1-a)/\vartheta>0$.  Monotone convergence first in $k$ and then in $T$
shows that the feedback cost is finite.  Finally, letting $k,T\to\infty$ in the equality and using
\eqref{eq:terminal-vanishes-verification} gives
$J_\lambda(X;\pi^*)=v_\lambda(X)$.
\end{proof}

\end{document}
